% ═══════════════════════════════════════════════════════════════
% ARBEITSBLATT: Atommodelle (Physik Klasse 10)
% Begleitend zum digitalen Lernpfad
% ═══════════════════════════════════════════════════════════════

\documentclass[10pt, a4paper]{article}

% ═══════════════════════════════════════════════════════════════
% PAKETE
% ═══════════════════════════════════════════════════════════════

\usepackage[utf8]{inputenc}
\usepackage[T1]{fontenc}
\usepackage[ngerman]{babel}
\usepackage{helvet}
\renewcommand{\familydefault}{\sfdefault}

\usepackage[margin=1.5cm, top=1.8cm, bottom=1.5cm]{geometry}
\usepackage{enumitem}
\usepackage{tabularx}
\usepackage{booktabs}
\usepackage{array}
\usepackage{multicol}

\usepackage{xcolor}
\usepackage{tcolorbox}
\usepackage{amssymb}
\usepackage{fancyhdr}
\usepackage{lastpage}
\usepackage[normalem]{ulem}

% Kompaktere Listen
\setlist{nosep, leftmargin=1.5em}

% ═══════════════════════════════════════════════════════════════
% FARBEN
% ═══════════════════════════════════════════════════════════════

\definecolor{physik}{HTML}{2c5aa0}
\definecolor{grau}{HTML}{888888}
\definecolor{hellgrau}{HTML}{f0f0f0}
\definecolor{success}{HTML}{2a7a4b}
\definecolor{fehler}{HTML}{c62828}
\definecolor{fehlerbg}{HTML}{ffebee}

\newcommand{\fachfarbe}{physik}

% ═══════════════════════════════════════════════════════════════
% VARIABLEN
% ═══════════════════════════════════════════════════════════════

\newcommand{\thema}{Atommodelle}
\newcommand{\fach}{Physik}
\newcommand{\klasse}{Klasse 10}
\newcommand{\stunde}{Stunde 1-2}

% ═══════════════════════════════════════════════════════════════
% KOPF- UND FUSSZEILE
% ═══════════════════════════════════════════════════════════════

\pagestyle{fancy}
\fancyhf{}
\renewcommand{\headrulewidth}{0.4pt}
\renewcommand{\footrulewidth}{0pt}

\lhead{\small\fach{} \klasse{} -- \thema}
\rhead{\small Name: \rule{4cm}{0.4pt}}

\cfoot{\small\thepage/\pageref{LastPage}}

% ═══════════════════════════════════════════════════════════════
% BEFEHLE
% ═══════════════════════════════════════════════════════════════

% Lücke
\newcommand{\luecke}[1]{\rule{#1}{0.4pt}}

% Physik-Highlight
\newcommand{\ph}[1]{\textcolor{\fachfarbe}{\textbf{#1}}}

% Merksatz-Box (zum Ausfüllen)
\newtcolorbox{merksatz}[1][]{
    colback=hellgrau,
    colframe=\fachfarbe,
    boxrule=0.8pt,
    arc=0pt,
    left=6pt, right=6pt, top=6pt, bottom=6pt,
    fonttitle=\bfseries\small,
    title=#1
}

% Fehler-Box
\newtcolorbox{fehlerbox}{
    colback=fehlerbg,
    colframe=fehler,
    boxrule=0.8pt,
    arc=0pt,
    left=6pt, right=6pt, top=4pt, bottom=4pt,
    fonttitle=\bfseries\small,
    title=Typische Fehler
}

% Abschnittstitel
\newcommand{\abschnitt}[1]{%
    \vspace{8pt}
    \noindent\textbf{\textcolor{\fachfarbe}{#1}}
    \vspace{4pt}
}

% ═══════════════════════════════════════════════════════════════
% DOKUMENT
% ═══════════════════════════════════════════════════════════════

\begin{document}

% ═══════════════════════════════════════════════════════════════
% TITEL
% ═══════════════════════════════════════════════════════════════

\begin{center}
{\large\bfseries Arbeitsblatt: \thema{} -- \stunde}\\[2pt]
{\small\textcolor{grau}{Begleitend zum Lernpfad -- in den Hefter übernehmen}}
\end{center}

\vspace{-2pt}
\hrule
\vspace{8pt}

% ═══════════════════════════════════════════════════════════════
% ZWEISPALTIG
% ═══════════════════════════════════════════════════════════════

\begin{multicols}{2}
\small

% ─────────────────────────────────────────────────────────────
% KASTEN 1: Die vier Atommodelle
% ─────────────────────────────────────────────────────────────

\abschnitt{Kasten 1: Die vier Atommodelle}

\begin{merksatz}[Historische Entwicklung]
Ergänze die Tabelle aus dem Lernpfad:

\vspace{6pt}

\renewcommand{\arraystretch}{1.8}
\begin{tabular}{|p{1.8cm}|p{0.8cm}|p{3.5cm}|}
\hline
\textbf{Modell} & \textbf{Jahr} & \textbf{Kernaussage} \\
\hline
Dalton & & \\
\hline
Thomson & & \\
\hline
Rutherford & & \\
\hline
Bohr & & \\
\hline
\end{tabular}
\renewcommand{\arraystretch}{1}
\end{merksatz}

\vspace{6pt}

% ─────────────────────────────────────────────────────────────
% KASTEN 2: Rutherford-Streuversuch
% ─────────────────────────────────────────────────────────────

\abschnitt{Kasten 2: Rutherford-Streuversuch}

\begin{merksatz}[Versuch und Ergebnisse]
\textbf{Aufbau:}\\
\luecke{2cm}-Teilchen werden auf eine dünne\\
\luecke{2cm}-Folie geschossen.

\vspace{6pt}

\textbf{Beobachtungen $\rightarrow$ Schlussfolgerungen:}

\vspace{4pt}

\renewcommand{\arraystretch}{1.6}
\begin{tabular}{|p{2.8cm}|p{3.2cm}|}
\hline
\textbf{Beobachtung} & \textbf{Schlussfolgerung} \\
\hline
Die meisten gehen durch & \\
\hline
Wenige werden stark abgelenkt & \\
\hline
\end{tabular}
\renewcommand{\arraystretch}{1}

\vspace{6pt}

\textbf{Größenverhältnis:}\\
Kerndurchmesser: $\approx$ \luecke{1.5cm} m\\
Atomdurchmesser: $\approx$ \luecke{1.5cm} m
\end{merksatz}

\vspace{6pt}

% ─────────────────────────────────────────────────────────────
% KASTEN 3: Kernaufbau
% ─────────────────────────────────────────────────────────────

\abschnitt{Kasten 3: Kernaufbau}

\begin{merksatz}[Bestandteile des Atoms]
\renewcommand{\arraystretch}{1.5}
\begin{tabular}{|l|c|c|c|}
\hline
\textbf{Teilchen} & \textbf{Ladung} & \textbf{Masse} & \textbf{Ort} \\
\hline
Proton & & & \\
\hline
Neutron & & & \\
\hline
Elektron & & & \\
\hline
\end{tabular}
\renewcommand{\arraystretch}{1}

\vspace{8pt}

\textbf{Schreibweise:} $^{A}_{Z}$X

\vspace{4pt}

A = \luecke{3cm} (Protonen + Neutronen)\\
Z = \luecke{3cm} (Anzahl Protonen)\\
N = \luecke{3cm} (Formel: \luecke{1.5cm})
\end{merksatz}

% ═══════════════════════════════════════════════════════════
% SPALTE 2
% ═══════════════════════════════════════════════════════════

\columnbreak

% ─────────────────────────────────────────────────────────────
% ÜBUNG 1: Zuordnung (AFB I)
% ─────────────────────────────────────────────────────────────

\abschnitt{Übung 1: Zeitstrahl ergänzen}

Ordne die Jahreszahlen und Modelle chronologisch:

\vspace{4pt}

\textit{Jahreszahlen: 1803, 1897, 1911, 1913}

\vspace{6pt}

\begin{center}
\luecke{1.2cm} $\rightarrow$ \luecke{1.2cm} $\rightarrow$ \luecke{1.2cm} $\rightarrow$ \luecke{1.2cm}

\vspace{4pt}

\luecke{1.5cm} $\rightarrow$ \luecke{1.5cm} $\rightarrow$ \luecke{1.5cm} $\rightarrow$ \luecke{1.5cm}
\end{center}

\vspace{6pt}

% ─────────────────────────────────────────────────────────────
% ÜBUNG 2: Lückentext (AFB II)
% ─────────────────────────────────────────────────────────────

\abschnitt{Übung 2: Rutherford erklären}

Ergänze den Lückentext:

\vspace{4pt}

Beim Streuversuch wurden \luecke{1.5cm}-Teilchen

auf eine dünne Goldfolie geschossen. Die meisten

gingen \luecke{2cm} durch. Das zeigt:

Das Atom ist größtenteils \luecke{2cm}.

Wenige Teilchen wurden stark \luecke{1.8cm}.

Das zeigt: Die positive Ladung ist in einem

\luecke{2cm} Kern konzentriert.

\vspace{8pt}

% ─────────────────────────────────────────────────────────────
% ÜBUNG 3: Kernaufbau (AFB II)
% ─────────────────────────────────────────────────────────────

\abschnitt{Übung 3: Teilchenzahlen bestimmen}

Bestimme für die folgenden Nuklide:

\vspace{6pt}

\begin{enumerate}[leftmargin=2em]
\item $^{12}_{6}$C: \quad p = \luecke{0.8cm} \quad n = \luecke{0.8cm} \quad e$^-$ = \luecke{0.8cm}

\vspace{2pt}

\item $^{235}_{92}$U: \quad p = \luecke{0.8cm} \quad n = \luecke{0.8cm} \quad e$^-$ = \luecke{0.8cm}

\vspace{2pt}

\item $^{56}_{26}$Fe: \quad p = \luecke{0.8cm} \quad n = \luecke{0.8cm} \quad e$^-$ = \luecke{0.8cm}
\end{enumerate}

\vspace{8pt}

% ─────────────────────────────────────────────────────────────
% ÜBUNG 4: Fehler begründen (AFB III)
% ─────────────────────────────────────────────────────────────

\abschnitt{Übung 4: Fehler erklären (AFB III)}

\begin{fehlerbox}
Erkläre, warum diese Aussagen falsch sind:

\vspace{6pt}

\begin{enumerate}[leftmargin=1.5em]
\item \textit{„Das Thomson-Modell wurde durch die Entdeckung des Elektrons widerlegt."}

\vspace{4pt}
\luecke{5.5cm}

\vspace{6pt}

\item \textit{„Im Rutherford-Versuch wurden die meisten Teilchen stark abgelenkt."}

\vspace{4pt}
\luecke{5.5cm}

\vspace{6pt}

\item \textit{„Die Massenzahl A gibt die Anzahl der Elektronen an."}

\vspace{4pt}
\luecke{5.5cm}
\end{enumerate}
\end{fehlerbox}

\vspace{8pt}

% ─────────────────────────────────────────────────────────────
% FREIES SCHREIBEN (AFB III)
% ─────────────────────────────────────────────────────────────

\abschnitt{Übung 5: Transfer}

Erkläre in 2-3 Sätzen: Warum musste nach dem Rutherford-Versuch ein neues Atommodell entwickelt werden?

\vspace{4pt}

\luecke{5.8cm}

\vspace{6pt}

\luecke{5.8cm}

\vspace{6pt}

\luecke{5.8cm}

\end{multicols}

% ═══════════════════════════════════════════════════════════════
% SEITE 2: SIMULATIONSPROTOKOLL
% ═══════════════════════════════════════════════════════════════

\newpage

\begin{center}
{\large\bfseries Simulationsprotokoll: Rutherford-Streuversuch}\\[2pt]
{\small\textcolor{grau}{PhET-Simulation: phet.colorado.edu/de/simulations/rutherford-scattering}}
\end{center}

\vspace{-2pt}
\hrule
\vspace{12pt}

% ─────────────────────────────────────────────────────────────
% VERSUCHSAUFBAU
% ─────────────────────────────────────────────────────────────

\abschnitt{Teil A: Versuchsaufbau verstehen}

\begin{merksatz}[Aufbau skizzieren]
Zeichne den Versuchsaufbau in die Box. Beschrifte: Strahlungsquelle, Alpha-Teilchen, Goldfolie, Zinksulfidschirm.

\vspace{80pt}

\end{merksatz}

\vspace{8pt}

% ─────────────────────────────────────────────────────────────
% VORHERSAGE (POE: Predict)
% ─────────────────────────────────────────────────────────────

\abschnitt{Teil B: Vorhersage (VOR der Simulation)}

\textbf{Wenn das Thomson-Modell stimmt} (gleichmäßig verteilte positive Ladung):

Was erwartest du? Kreuze an:

\vspace{4pt}

$\square$ Alle Teilchen werden stark abgelenkt

$\square$ Alle Teilchen gehen unabgelenkt durch

$\square$ Alle Teilchen werden gleichmäßig leicht abgelenkt

\vspace{6pt}

Begründe deine Vorhersage: \luecke{10cm}

\vspace{8pt}

% ─────────────────────────────────────────────────────────────
% BEOBACHTUNG (POE: Observe)
% ─────────────────────────────────────────────────────────────

\abschnitt{Teil C: Beobachtung (WÄHREND der Simulation)}

Starte die Simulation und beobachte. Zähle bei 100 abgefeuerten Teilchen:

\vspace{6pt}

\begin{tabular}{|p{6cm}|c|c|}
\hline
\textbf{Beobachtung} & \textbf{Anzahl} & \textbf{Anteil} \\
\hline
Teilchen gehen unabgelenkt durch & \luecke{1.5cm} & \luecke{1.5cm} \% \\
\hline
Teilchen werden leicht abgelenkt (<45°) & \luecke{1.5cm} & \luecke{1.5cm} \% \\
\hline
Teilchen werden stark abgelenkt (>45°) & \luecke{1.5cm} & \luecke{1.5cm} \% \\
\hline
Teilchen werden zurückgeworfen (>90°) & \luecke{1.5cm} & \luecke{1.5cm} \% \\
\hline
\end{tabular}

\vspace{8pt}

Stimmt deine Vorhersage mit der Beobachtung überein? $\square$ Ja \quad $\square$ Nein

\vspace{8pt}

% ─────────────────────────────────────────────────────────────
% ERKLÄRUNG (POE: Explain)
% ─────────────────────────────────────────────────────────────

\abschnitt{Teil D: Erklärung (NACH der Simulation)}

\textbf{1. Warum gehen die meisten Teilchen unabgelenkt durch?}

\vspace{4pt}
\luecke{14cm}

\vspace{8pt}

\textbf{2. Warum werden manche Teilchen stark abgelenkt oder zurückgeworfen?}

\vspace{4pt}
\luecke{14cm}

\vspace{8pt}

\textbf{3. Was bedeutet das für den Aufbau des Atoms?}

\vspace{4pt}
\luecke{14cm}

\vspace{12pt}

% ─────────────────────────────────────────────────────────────
% ZUSATZEXPERIMENT
% ─────────────────────────────────────────────────────────────

\abschnitt{Teil E: Zusatzexperiment (für Schnelle)}

Verändere in der Simulation die \textbf{Anzahl der Protonen} im Kern.

\vspace{6pt}

Was passiert bei \textbf{mehr} Protonen? \luecke{8cm}

\vspace{6pt}

Erkläre, warum: \luecke{10cm}

\end{document}
